\documentclass{leveldoc}
\DocUses{foundations,first_bend,slingshot}
\DocFingerprint{ef7f8d0f0a55d414}

\begin{document}

\section{What the level asks}\label{sec:asks}

The game's text:
\leveldescription

One ion is launched from $(\val{xZero}, \val{yZero})$ along $+x$ with $T_0 = \val{Tzero}$, in
Newton's mechanics ($c = \infty$). It must enter the detector
$\val{xDetLo} \le x \le \val{xDetHi}$, $\val{yDetLo} \le y \le \val{yDetHi}$ under both of the
level's disturbances: ``\val{distOne}'', no outside field, and ``\val{distTwo}'', a uniform field
$\vect{E} = (0, \val{Ey})$ across the beam line (\physref{2.3}: every shot is flown under every
disturbance, and the level is solved only if every flight arrives). The player places up to
three charges from the listed magnitudes. A flight ends when the ion enters the detector,
touches a charge, or leaves the world (\xsecref{foundations}{sec:board}).

The reference places one positive charge $Q = \val{Q}$ at $(\val{xC}, \val{yC})$, below and
behind the launch: its repulsion lifts the ion enough that it enters the detector with the field
off, near the detector's middle, and still enters it, lower, with the field pushing down. With
the field off this is Level 1's flight about one charge, in closed form
(\levelref{first_bend}); with it on, a numerical march.

\section{Setup}\label{sec:setup}

The level file, field by field (the appendix reproduces it):
\setuptable

The ion has charge $q = \val{q}$, mass $m = \val{m}$ and radius $0$; its momentum is
$p_0 = \sqrt{2mT_0} = \val{pZero}$ and its speed $v_0 = \val{vZero}$. The coupling with the
reference charge is $\kappa = qQ = \val{kappa}$ (repulsive); the charge's radius is $\val{R}$.
The field, when on, gives the ion the acceleration $qE_y/m = \val{accel}$. The world is
$-1 \le x \le \val{xWorld}$, $-1 \le y \le \val{yWorld}$.

\begin{figure}[htbp]
\centering
\begin{tikzpicture}
\begin{axis}[axis equal image, width=0.92\linewidth, xmin=-1.6, xmax=31.6, ymin=-1.6, ymax=21.6,
  xtick={0,5,...,30}, ytick={0,5,...,20}, tick label style={font=\small},
  xlabel={$x$}, ylabel={$y$}, label style={font=\small}]
  \draw[black!8] (0,0) grid (30,20);
  \draw[black!45, dashed] (-1,-1) rectangle (\rawval{xWorld},\rawval{yWorld});
  \fill[green!50!black, fill opacity=0.18] (\rawval{xDetLo},\rawval{yDetLo}) rectangle (\rawval{xDetHi},\rawval{yDetHi});
  \draw[green!40!black] (\rawval{xDetLo},\rawval{yDetLo}) rectangle (\rawval{xDetHi},\rawval{yDetHi});
  \addplot[thick, blue!70!black] table {pathoff.dat};
  \addplot[thick, red!70!black] table {pathon.dat};
  \fill[red!75!black] (\rawval{xC},\rawval{yC}) circle[radius=\rawval{R}];
  \fill (\rawval{xZero},\rawval{yZero}) circle[radius=0.15];
\end{axis}
\end{tikzpicture}
\caption{The reference charge (red disc) and the two flights: with the field off (blue), a
Kepler hyperbola about the charge; with it on (red), pushed lower, still into the detector.}
\label{fig:board}
\end{figure}

\section{Principles and model}\label{sec:model}

\begin{principles}
  \item Newton's equation of motion with the force $q(\vect{E}_Q + \vect{E}_\mathrm{out})$, the
    charge's field \xeqref{foundations}{eq:coulomb} and the uniform outside field (\physref{2.3}).
  \item The orbit about one charge, \xeqref{first_bend}{eq:general} in its Newtonian limit, with
    the time \xeqref{first_bend}{eq:dtdth} and the momentum \xeqref{first_bend}{eq:momentum}.
  \item The board's rules; the classical Runge--Kutta method, cubic Hermite interpolation, the
    error estimate from the step $2h$ and Simpson's rule (\xsecref{foundations}{sec:board},
    \xsecref{foundations}{sec:numerics}), as in \xsecref{slingshot}{sec:method}.
\end{principles}

The charge is held in place, the ion is a point, and the outside field is uniform and static
under each disturbance: an exact solution of Maxwell's equations in vacuum, so it adds no
approximation (\physref{2.3}).

\section{Derivation}\label{sec:derivation}

\subsection{The size of the disturbance}\label{sec:size}

Crossing to the detector at about $v_0$ takes $t \approx \val{tCross}$; a uniform push
$a = \val{accel}$ alone moves the ion by $\tfrac12 a t^2 = \val{dropEst}$ in that time, comparable
to the detector's height: the setup must work at two heights. On the empty board the field-on
flight is exactly such a parabola, $y = y_0 + \tfrac12 a (x/v_0)^2$, which leaves the world's right
edge at $y = \val{bareYOn}$ (as the table's empty-board row says).

\subsection{The field off: Level 1's closed form}\label{sec:kepler}

With one charge and no outside field, \xeqref{first_bend}{eq:orbit} in Newton's limit
($\Gamma = 1$, $W/c^2 \to m$) is Binet's equation $u'' + u = -\kappa m/L^2$, $u = 1/r$, with the
angular momentum $L = \val{L}$ about the charge; its solution from the launch values is
\xeqref{first_bend}{eq:general}: $u(\theta) = B + (u_0 - B)\cos\theta + u'_0\sin\theta$ with
$B = -\kappa m/L^2 = \val{Bk}$. Stepping $\theta$ and bisecting, the orbit enters the detector after
$\theta = \val{thK}$ degrees, at $y = \val{yK}$; Simpson's rule on
$\dd t/\dd\theta = m/(L u^2)$ gives $t = \val{tK}$ (its error $\val{tKerr}$).

\subsection{The field on: the march}\label{sec:march}

With the field on, the energy
\begin{equation}\label{eq:energy-stray}
  W = \frac{p^2}{2m} + \frac{\kappa}{\abs{\vect{x} - \vect{r}_Q}} - qE_y\,y
\end{equation}
is conserved, but the angular momentum about the charge is not, and the orbit has no form as
simple as Binet's. Both flights are marched by \xeqref{foundations}{eq:rk4} at $h = \val{H}$,
the entry located on the cubic Hermite interpolant; the error is estimated against the step
$2h$. With the field off the march agrees with the closed form to $\val{kDiff}$.

\section{The reference solution by hand}\label{sec:hand}

\begin{procedure}
  \item\label{st:off} \emph{Field off.} The ion enters the detector through its \val{faceOff}
    face at $(\val{xOff}, \val{yOff})$ at $t = \val{tOff}$ (step estimate $\val{errOff}$; the
    energy drifts by $\val{wDevOff}$), as the closed form says.
  \item\label{st:on} \emph{Field on.} The ion enters through the \val{faceOn} face at
    $(\val{xOn}, \val{yOn})$ at $t = \val{tOn}$ (step estimate $\val{errOn}$; the energy
    \eqref{eq:energy-stray} drifts by $\val{wDevOn}$): lower, inside the detector still.
\end{procedure}

\begin{result}[The reference flights]
Under both disturbances the ion enters the detector: at $(\val{xOff}, \val{yOff})$ with the field
off and at $(\val{xOn}, \val{yOn})$ with it on. The level is solved.
\end{result}

\section{Comparison with the game}\label{sec:comparison}

\comparisontable

All differences are below $10^{-10}$, the game's verified accuracy (Foundations §6); our
estimates are a few $10^{-13}$. The tolerance, $10^{-9}$, is ten times the larger of the two
computations' estimates, fixed before comparing. The largest difference, $\val{largest}$, is the
game's: its verification controls each step's error to
$10^{-12}$ (the level's \code{tolerances.verify}), and over the flight's few dozen steps and the
location of the entry those errors add up to some $10^{-11}$, while our march's estimate is a
hundred times smaller.

\section{Other solutions}\label{sec:others}

On the empty board the ion flies straight with the field off, leaving the right edge at
$y = \val{bareYOff}$, and falls on the parabola of \S\ref{sec:size} with it on: neither enters the
detector. A solution must lift the ion by about the detector's height above the launch line, so
that the field-off flight enters in the detector's upper part and the field-on one, about
$\val{dropEst}$ lower by \S\ref{sec:size}, still above its lower edge. One charge does it in two
ways, by the same march, checked against the game in the comparison table:
\begin{itemize}
  \item a positive charge below the line pushes the ion up: the reference, or a small one,
    $Q = \val{QPush}$ at $(\val{xPush}, \val{yPush})$ just below the launch, which enters at
    $y = \val{yPushOff}$ with the field off and $\val{yPushOn}$ with it on;
  \item a negative charge above the line pulls it up: $Q = \val{QPull}$ at
    $(\val{xPull}, \val{yPull})$ enters at $y = \val{yPullOff}$ and $\val{yPullOn}$.
\end{itemize}
In both, a weaker charge nearer the line does what a stronger one farther away does: a charge
at distance $b$ from the line, level with the launch, gives the ion flying past it at $v_0$ the
sideways impulse (the transverse Coulomb force integrated along the undeflected line)
\begin{equation}\label{eq:impulse}
  \Delta p_\perp = \int_0^\infty \frac{\kappa b\,\dd t}{(b^2 + v_0^2 t^2)^{3/2}} = \frac{\kappa}{b v_0},
\end{equation}
so the lift asks for about the same $Q/b$. A map of all solutions
would take the game's verdicts, which this document does not reproduce.

\section{What the game computes}\label{sec:engine}

The game adds the disturbance's uniform field to the level's sources (\physref{2.3}) and
integrates Newton's equations with its adaptive eighth-order Runge--Kutta method (\physref{5},
\physref{6}, \physref{7}), each disturbance's flight verified on its own.

\section{Further reading}\label{sec:reading}

The Kepler problem and Binet's equation: \Goldstein{§3.5}, \levelref{first_bend}.

\levelfileappendix

\end{document}
